\chapter{Πρόβλεψη πωλήσεων και οικονομική ανάλυση}

% TODO(PR9, PR10): Συγγραφή μαζί με την υλοποίηση της πρόβλεψης και των
% οικονομικών σεναρίων.

\section{Πρόβλεψη χρονοσειρών πωλήσεων}
\label{sec:forecasting}

Η χρονοσειρά που προβλέπεται είναι η καθαρή και όχι η ακαθάριστη αξία
πωλήσεων. Οι επιστροφές αποτελούν χρήματα που επεστράφησαν, και μια πρόβλεψη
επί της ακαθάριστης αξίας θα υπερεκτιμούσε συστηματικά το ποσό που η
επιχείρηση μπορεί να αναμένει ότι θα διατηρήσει.

Εφαρμόζεται εκθετική εξομάλυνση κατά Holt-Winters, με προσθετική τάση και,
όπου το ιστορικό επαρκεί, προσθετική εποχικότητα. Ο έλεγχος επάρκειας είναι
ρητός: μια εποχική προσαρμογή απαιτεί τουλάχιστον δύο πλήρεις κύκλους, και
όταν αυτοί δεν υπάρχουν το μοντέλο υποχωρεί σε προσαρμογή μόνο με τάση αντί να
αποτύχει.

\section{Αξιολόγηση της πρόβλεψης}

Η αξιολόγηση γίνεται εκτός δείγματος. Ένα τμήμα στο τέλος της χρονοσειράς
δεσμεύεται, το μοντέλο προσαρμόζεται αποκλειστικά στα προηγούμενα σημεία, και
τα σφάλματα υπολογίζονται επί των δεσμευμένων παρατηρήσεων. Η τελική πρόβλεψη
επαναπροσαρμόζεται στο σύνολο της σειράς, καθώς δεν υπάρχει λόγος απόρριψης των
πλέον πρόσφατων παρατηρήσεων κατά την πρόβλεψη του μέλλοντος.

Αναφέρονται τρία μέτρα σφάλματος: το μέσο απόλυτο σφάλμα
(\textenglish{MAE}), η ρίζα του μέσου τετραγωνικού σφάλματος
(\textenglish{RMSE}), η οποία τιμωρεί βαρύτερα τις μεγάλες αποκλίσεις, και το
μέσο απόλυτο ποσοστιαίο σφάλμα (\textenglish{MAPE}).

Το \textenglish{MAPE} δεν ορίζεται όταν κάποια πραγματική τιμή ισούται με
μηδέν. Στην περίπτωση αυτή αναφέρεται ρητά ως μη ορισμένο. Η αντικατάστασή του
με μικρή σταθερά ή η σιωπηρή παράλειψη των αντίστοιχων σημείων θα κολάκευε
τεχνηέντως το μοντέλο.

\subsection{Σύγκριση με αφελή βασική γραμμή}

Κάθε πρόβλεψη συνοδεύεται από αφελή βασική γραμμή, η οποία απλώς επαναλαμβάνει
την τελευταία παρατηρούμενη τιμή. Η επιλογή αυτή είναι μεθοδολογικά κρίσιμη:
ένα μοντέλο που δεν υπερβαίνει την πρόταση «το αύριο μοιάζει με το σήμερα» δεν
έχει αποδείξει τίποτε, και ένα \textenglish{MAPE} που παρατίθεται χωρίς τέτοια
σύγκριση αποτελεί μη διαψεύσιμο αριθμό.

Η βελτίωση εκφράζεται μέσω του δείκτη δεξιότητας
(\textenglish{skill score}), ως το κλασματικό κέρδος στο μέσο απόλυτο σφάλμα
έναντι της βασικής γραμμής. Ο δείκτης μπορεί να λάβει αρνητική τιμή, και στην
περίπτωση αυτή αναφέρεται ως έχει.

\section{Μετρήσιμοι οικονομικοί δείκτες}
\label{sec:measured-economics}

Οι δείκτες της ενότητας αυτής υπολογίζονται απευθείας από τα δεδομένα
συναλλαγών. Καμία παραδοχή δεν εισάγεται σε αυτό το σημείο, και η διάκριση
από την ενότητα~\ref{sec:profitability-scenarios} που ακολουθεί είναι
ουσιώδης.

\subsection{Ακαθάριστη αξία πωλήσεων}

Η ακαθάριστη αξία πωλήσεων (\textenglish{Gross Merchandise Value, GMV})
αναφέρεται πριν από τις επιστροφές, σύμφωνα με τον καθιερωμένο ορισμό. Η αξία
των επιστροφών και η καθαρή αξία παρατίθενται δίπλα της, ώστε κανένα από τα
δύο μεγέθη να μην αποκρύπτεται. Η επιλογή αυτή είναι συνεπής με τη σχεδίαση
του σχήματος δεδομένων, όπου οι επιστροφές αποτελούν αυτοτελείς εγγραφές και
δεν αφαιρούνται σιωπηρά από τα έσοδα.

\subsection{Ποσοστό επαναληπτικών αγορών}

Ως επαναλαμβανόμενος θεωρείται ο πελάτης με περισσότερες από μία παραγγελίες.
Στον υπολογισμό συμμετέχουν μόνο οι αναγνωρισμένοι πελάτες: οι ανώνυμες
πωλήσεις δεν μπορούν να αποδοθούν σε πρόσωπο, και η συμπερίληψή τους θα
υποεκτιμούσε συστηματικά την επαναληπτικότητα.

\subsection{Διατήρηση πελατών ανά κοόρτη}

Οι πελάτες ομαδοποιούνται με βάση τον μήνα της πρώτης τους αγοράς και
παρακολουθούνται στους επόμενους μήνες. Ο δείκτης περιόδου μετρά τους
συμπληρωμένους μήνες από την απόκτηση, οπότε η τιμή μηδέν αντιστοιχεί στον
ίδιο τον μήνα της κοόρτης. Το πλέγμα που προκύπτει είναι εξ ορισμού
τριγωνικό, καθώς οι νεότερες κοόρτες διαθέτουν λιγότερους μήνες
παρακολούθησης.

\subsection{Ανάλυση RFM}

Για κάθε πελάτη υπολογίζονται η προσφατότητα (\textenglish{recency}), η
συχνότητα (\textenglish{frequency}) και η χρηματική αξία
(\textenglish{monetary}). Η προσφατότητα μετράται σε συμπληρωμένες ημέρες από
την τελευταία αγορά, ως προς ημερομηνία αναφοράς που εξ ορισμού είναι η
τελευταία παραγγελία του συνόλου δεδομένων· η επιλογή αυτή διατηρεί τα
αποτελέσματα αναπαραγώγιμα για ένα ιστορικό σύνολο δεδομένων, αντί να
μεταβάλλονται με την τρέχουσα ημερομηνία.

Η χρηματική αξία αποτελεί την παρατηρούμενη αξία του πελάτη μέχρι σήμερα και
όχι πρόβλεψη αξίας διάρκειας ζωής. Δεν εμπλέκεται κανένα μοντέλο σε αυτό το
στάδιο. Τα ίδια χαρακτηριστικά τροφοδοτούν τη συσταδοποίηση δεύτερου επιπέδου
του κεφαλαίου~\ref{chap:clustering}, ακολουθώντας την προσέγγιση των
\textcite{chen2012data}.

\section{Ανάλυση κερδοφορίας βάσει σεναρίων}
\label{sec:profitability-scenarios}

Τα ερωτήματα «είναι η επιχείρηση κερδοφόρα» και «πότε θα γίνει κερδοφόρα»
\emph{δεν} απαντώνται από τα διαθέσιμα δεδομένα. Το σύνολο δεδομένων καταγράφει
τι πωλήθηκε και σε ποια τιμή, δεν καταγράφει όμως ούτε κόστος πωληθέντων ούτε
δαπάνες marketing. Οποιοδήποτε μέγεθος κερδοφορίας παρουσιαζόταν ως μετρημένο
θα αποτελούσε επινόηση με το ένδυμα ευρήματος.

Αυτό που μπορεί να γίνει με ακαδημαϊκή εντιμότητα είναι ανάλυση ευαισθησίας:
δήλωση ρητών παραδοχών, μεταβολή τους, και καταγραφή του τι προκύπτει
\emph{εφόσον} αυτές ισχύουν.

\subsection{Δηλωμένες παραδοχές}

Οι παραδοχές είναι δύο, το μεικτό περιθώριο κέρδους και το κόστος απόκτησης
πελάτη, και προαιρετικά ένα σταθερό λειτουργικό κόστος ανά περίοδο. Δεν
προκύπτουν από τα δεδομένα και δεν υπονοούνται πουθενά: παρέχονται από τον
χρήστη και επιστρέφονται μαζί με κάθε αποτέλεσμα, ώστε παραδοχή και συμπέρασμα
να μην μπορούν να αποσυνδεθούν.

Κάθε απόκριση του συστήματος φέρει επιπλέον ρητή δήλωση ότι πρόκειται για
σενάριο και όχι για μέτρηση. Η δήλωση αυτή ταξιδεύει μαζί με τα δεδομένα σε
όλα τα επίπεδα, από τον κώδικα έως το API, ακριβώς για να μην μπορεί να
αφαιρεθεί εκ παραδρομής.

\subsection{Σημεία αναφοράς της αγοράς}

Ως πλαίσιο σύγκρισης χρησιμοποιούνται δημοσιευμένες εκτιμήσεις της αγοράς: το
μεικτό περιθώριο κυμαίνεται περίπου από 15-30\% στα ηλεκτρονικά έως 65-85\%
στα καλλυντικά και τα συμπληρώματα, ο λόγος αξίας διάρκειας ζωής προς κόστος
απόκτησης θεωρείται υγιής περί το 3:1, και η περίοδος απόσβεσης του κόστους
απόκτησης τυπικά κινείται μεταξύ 6 και 12 μηνών.\footnote{Εκτιμήσεις κλάδου από
επαγγελματικές πηγές, οι οποίες δεν αποτελούν αξιολογημένη από ομοτίμους
έρευνα. Παρατίθενται ως πλαίσιο για την οριοθέτηση εύλογων τιμών των
παραδοχών και όχι ως μεγέθη της συγκεκριμένης επιχείρησης.}

\subsection{Αξία διάρκειας ζωής και απόσβεση}

Η αξία διάρκειας ζωής υπολογίζεται επί του μεικτού κέρδους και όχι επί των
εσόδων. Η σύγκριση εσόδων με το κόστος απόκτησης θα υπερεκτιμούσε τον λόγο
ακριβώς κατά το κόστος πωληθέντων, και αποτελεί τον συνηθέστερο τρόπο με τον
οποίο ο δείκτης αυτός εμφανίζεται τεχνηέντως ευνοϊκός.

Η περίοδος απόσβεσης προϋποθέτει ρυθμό κέρδους, επομένως απαιτεί και το μέσο
πλήθος περιόδων δραστηριότητας του πελάτη. Όταν αυτό δεν είναι διαθέσιμο, η
περίοδος απόσβεσης αναφέρεται ως μη ορισμένη αντί να εκτιμηθεί αυθαίρετα.

\subsection{Πλέγμα ευαισθησίας}

Ένα μεμονωμένο σενάριο προσκαλεί την ανάγνωσή του ως πρόβλεψης. Για τον λόγο
αυτό τα αποτελέσματα παρουσιάζονται ως πλέγμα πάνω σε εύρος παραδοχών.

Το πλέγμα καθιστά ορατή την εξάρτηση του συμπεράσματος από τις παραδοχές: τα
ίδια ακριβώς παρατηρούμενα δεδομένα οδηγούν σε αντίθετα συμπεράσματα ως προς
το σημείο αναφοράς 3:1, ανάλογα και μόνο με το υποτιθέμενο περιθώριο κέρδους.
Η ιδιότητα αυτή επαληθεύεται αυτόματα στη σουίτα ελέγχων και αποτελεί το
κεντρικό συμπέρασμα της ενότητας.

\subsection{Προβολή στον χρόνο}

Η προβολή κερδοφορίας εκτείνεται στις καταγεγραμμένες και στις προβλεπόμενες
περιόδους ως ενιαία σειρά, με ρητή σήμανση του ορίου μεταξύ τους. Μια προβολή
που στηρίζεται σε προβλεπόμενα έσοδα είναι ασθενέστερη από μία που στηρίζεται
σε καταγεγραμμένα, και ο αναγνώστης δικαιούται να γνωρίζει πού βρίσκεται αυτό
το όριο. Για τον ίδιο λόγο αναφέρεται και ο δείκτης δεξιότητας της υποκείμενης
πρόβλεψης της ενότητας~\ref{sec:forecasting}.
